\documentclass[a4paper,12pt]{article}
\usepackage[T2A]{fontenc}
\usepackage[utf8]{inputenc}
\usepackage[russian]{babel}
\usepackage{amsmath,amssymb}
\renewcommand{\familydefault}{\sfdefault}
\usepackage{icomma}
\usepackage{geometry}
\geometry{top=22mm,bottom=18mm,left=20mm,right=20mm}
\setlength{\headheight}{25pt}
\usepackage{microtype}
\usepackage{xcolor}
\usepackage{tikz}
\usetikzlibrary{calc,arrows.meta,positioning,shapes.geometric}
\usepackage{tabularx,booktabs}
\usepackage{enumitem}
\usepackage{parskip}
\usepackage{fancyhdr}
\usepackage{setspace}

\definecolor{accent}{HTML}{2F6F6D}
\definecolor{darkaccent}{HTML}{244F4D}
\definecolor{textgray}{HTML}{303638}
\definecolor{softgray}{HTML}{EEF2F1}

\onehalfspacing
\setlength{\parindent}{0pt}
\setlist[itemize]{leftmargin=6mm,itemsep=2pt,topsep=3pt}
\setlist[enumerate]{leftmargin=8mm,itemsep=5pt,topsep=4pt}
\pagestyle{fancy}
\fancyhf{}
\fancyhead[L]{\small\color{textgray}Физика: движение тел}
\fancyhead[R]{\small\color{textgray}30.09.26}
\fancyfoot[C]{\small\thepage}
\renewcommand{\headrulewidth}{0.3pt}
\renewcommand{\headrule}{\hbox to\headwidth{\color{accent}\leaders\hrule height \headrulewidth\hfill}}

\newcommand{\sectionline}{\par\vspace{2mm}{\color{accent}\hrule height 0.7pt}\vspace{3mm}}
\newcommand{\checkitem}[1]{\item[$\square$] #1}
\newcommand{\unit}[1]{\,\text{#1}}

\tikzset{
  startstop/.style={ellipse, draw=accent, very thick, align=center, minimum width=38mm, minimum height=11mm, fill=softgray},
  process/.style={rectangle, rounded corners=3pt, draw=accent, very thick, align=center, text width=112mm, minimum height=12mm, fill=white},
  decision/.style={diamond, aspect=2.5, draw=accent, very thick, align=center, text width=42mm, inner sep=2pt, fill=softgray},
  arrow/.style={-{Stealth[length=3mm]}, very thick, draw=darkaccent},
  lab/.style={font=\small, text=darkaccent, fill=white, inner sep=1pt}
}

\begin{document}
\pagenumbering{gobble}

\begin{titlepage}
\thispagestyle{empty}
\centering
\vspace*{24mm}
{\color{accent}\Large ФИЗИКА\par}
\vspace{8mm}
{\Huge\bfseries Чек-лист\par}
\vspace{4mm}
{\LARGE\bfseries Движение протяжённых тел и условие встречи\par}
\vspace{14mm}
{\large Володя Хачатурян\par}
\vspace{5mm}
{\large 30 сентября 2026 года\par}
\vfill
{\small\color{textgray}Лёвин Артём Александрович эксклюзивно для Володи Хачатуряна\par}
\vspace{15mm}
\begin{tikzpicture}[remember picture,overlay]
  \draw[accent, line width=1.2pt]
    ($(current page.south west)+(18mm,18mm)$)
    .. controls ($(current page.south west)+(62mm,35mm)$)
    and ($(current page.south east)+(-60mm,5mm)$)
    .. ($(current page.south east)+(-18mm,20mm)$);
\end{tikzpicture}
\end{titlepage}
\clearpage
\pagenumbering{arabic}

\section*{1. Что нужно знать наизусть}
\sectionline
\begin{itemize}
\checkitem{Для равномерного движения координата: \(x=x_0+v_x t\). Знак \(v_x\) определяется направлением относительно выбранной оси \(Ox\).}
\checkitem{Для равноускоренного движения координата: \(x=x_0+v_{0x}t+\dfrac{a_x t^2}{2}\).}
\checkitem{Условие встречи двух тел: \(\boxed{x_1(t)=x_2(t)}\).}
\checkitem{При движении навстречу относительная скорость равна \(v_{\text{отн}}=v_1+v_2\).}
\checkitem{При движении в одном направлении относительная скорость по модулю равна \(v_{\text{отн}}=|v_1-v_2|\).}
\checkitem{В задачах с протяжёнными телами учитываются их длины. Для полного прохождения одного объекта мимо другого относительное перемещение обычно равно сумме их длин.}
\checkitem{Перед подстановкой все длины, скорости и время должны быть выражены в согласованных единицах.}
\end{itemize}

\section*{2. Движение протяжённых тел}
\sectionline
Пусть два протяжённых объекта длиной \(L_1\) и \(L_2\) проходят друг мимо друга. Если в условии нет дополнительных расстояний, рабочая модель имеет вид
\[
\boxed{t=\frac{L_1+L_2}{v_{\text{отн}}}}.
\]
Если один объект считается точкой (столб, человек), его длина принимается равной нулю. Если объект неподвижен, его скорость равна нулю.

Если \(d_{\text{до}}\) --- начальный зазор между передней точкой догоняющего и задней точкой впереди идущего объекта, а \(d_{\text{после}}\) --- требуемый конечный зазор между задней точкой обогнавшего и передней точкой второго объекта, то
\[
\boxed{t=\frac{L_1+L_2+d_{\text{до}}+d_{\text{после}}}{v_{\text{отн}}}}.
\]

\subsection*{Пример 1. Поезд проходит мимо столба}
Поезд движется со скоростью \(60\unit{км/ч}\) и проходит мимо столба за \(30\unit{с}\). Так как \(60\unit{км/ч}=\dfrac{50}{3}\unit{м/с}\),
\[
L=vt=\frac{50}{3}\cdot 30=500\unit{м}.
\]

\subsection*{Пример 2. Поезд проходит мимо лесополосы}
Скорость поезда \(90\unit{км/ч}=25\unit{м/с}\), длина лесополосы \(800\unit{м}\), время \(60\unit{с}\). За это время передняя точка поезда проходит
\[
S=25\cdot60=1500\unit{м}.
\]
Следовательно, длина поезда
\[
L=1500-800=700\unit{м}.
\]

\subsection*{Типичные ошибки}
\begin{itemize}
\item складывать скорости при движении в одном направлении;
\item забывать длину второго объекта;
\item использовать \(\unit{км/ч}\) вместе с секундами или метры вместе с часами без перевода;
\item путать путь обычной материальной точки и относительное перемещение, необходимое для полного обгона.
\end{itemize}

\clearpage
\thispagestyle{fancy}
\begin{center}
{\Large\bfseries\color{accent} Алгоритм: задача на движение протяжённых тел}
\end{center}
\vspace{7mm}
\begin{tikzpicture}[node distance=10mm, every node/.style={font=\small}]
\node[startstop] (start) {Прочитайте условие};
\node[process, below=of start] (len) {Определите длины \(L_1,L_2\) и дополнительные расстояния \(d_{\text{до}},d_{\text{после}}\), если они есть};
\node[decision, below=14mm of len] (dir) {Как движутся объекты?};
\node[process, below left=16mm and 10mm of dir, text width=50mm] (opp) {Навстречу:\\ \(v_{\text{отн}}=v_1+v_2\)};
\node[process, below right=16mm and 10mm of dir, text width=50mm] (same) {В одном направлении:\\ \(v_{\text{отн}}=|v_1-v_2|\)};
\node[process, below=24mm of dir] (units) {Приведите длины, скорости и время к согласованным единицам};
\node[process, below=of units] (dist) {Составьте относительное расстояние:\\ \(S_{\text{отн}}=L_1+L_2+d_{\text{до}}+d_{\text{после}}\)};
\node[process, below=of dist] (eq) {Используйте \(S_{\text{отн}}=v_{\text{отн}}t\) и найдите неизвестную величину};
\node[startstop, below=of eq] (check) {Проверьте единицы и физический смысл ответа};
\draw[arrow] (start)--(len);
\draw[arrow] (len)--(dir);
\draw[arrow] (dir)--node[lab,left]{навстречу}(opp);
\draw[arrow] (dir)--node[lab,right]{одно направление}(same);
\draw[arrow] (opp.south) |- (units.west);
\draw[arrow] (same.south) |- (units.east);
\draw[arrow] (units)--(dist);
\draw[arrow] (dist)--(eq);
\draw[arrow] (eq)--(check);
\end{tikzpicture}
\clearpage

\section*{3. Встреча двух тел через координаты}
\sectionline
Главная идея: два тела встречаются в тот момент, когда находятся в одной и той же точке пространства. Поэтому после выбора оси \(Ox\) нужно записать два уравнения координат и приравнять их.

Для равномерного движения:
\[
x_1=x_{01}+v_{1x}t,\qquad x_2=x_{02}+v_{2x}t,
\]
а условие встречи:
\[
\boxed{x_{01}+v_{1x}t=x_{02}+v_{2x}t}.
\]

\subsection*{Пример 3. Два автомобиля движутся навстречу}
Пункты \(A\) и \(B\) находятся на расстоянии \(120\unit{км}\). Из \(A\) автомобиль движется со скоростью \(90\unit{км/ч}\), из \(B\) навстречу ему --- со скоростью \(110\unit{км/ч}\). Направим ось \(Ox\) от \(A\) к \(B\):
\[
x_1=90t,\qquad x_2=120-110t.
\]
При встрече
\[
90t=120-110t,\qquad 200t=120,\qquad t=0{,}6\unit{ч}.
\]
Место встречи:
\[
x=90\cdot0{,}6=54\unit{км}
\]
от пункта \(A\).

\subsection*{Что особенно важно}
\begin{itemize}
\item Начальные координаты зависят от выбора начала отсчёта.
\item Проекция скорости положительна, если тело движется вдоль положительного направления оси, и отрицательна при движении против него.
\item Равенство координат работает и при встречном, и при попутном движении: направление уже учтено знаками скоростей.
\item После нахождения \(t\) место встречи можно получить подстановкой времени в любое из двух уравнений координаты.
\end{itemize}

\clearpage
\thispagestyle{fancy}
\begin{center}
{\Large\bfseries\color{accent} Алгоритм: найти время и место встречи}
\end{center}
\vspace{9mm}
\begin{tikzpicture}[node distance=11mm, every node/.style={font=\small}]
\node[startstop] (start2) {Изобразите ситуацию и выберите ось \(Ox\)};
\node[process, below=of start2] (x0) {Определите начальные координаты \(x_{01}\) и \(x_{02}\)};
\node[process, below=of x0] (sign) {Определите знаки проекций скоростей \(v_{1x}\) и \(v_{2x}\)};
\node[process, below=of sign] (eqs) {Запишите \(x_1(t)\) и \(x_2(t)\)};
\node[process, below=of eqs] (meet) {Используйте условие встречи \(x_1(t)=x_2(t)\)};
\node[process, below=of meet] (time) {Решите уравнение и найдите \(t\)};
\node[process, below=of time] (place) {Подставьте \(t\) в \(x_1(t)\) или \(x_2(t)\) и найдите координату встречи};
\node[startstop, below=of place] (check2) {Проверьте знак, единицы и положение точки встречи};
\draw[arrow] (start2)--(x0);
\draw[arrow] (x0)--(sign);
\draw[arrow] (sign)--(eqs);
\draw[arrow] (eqs)--(meet);
\draw[arrow] (meet)--(time);
\draw[arrow] (time)--(place);
\draw[arrow] (place)--(check2);
\end{tikzpicture}
\clearpage

\section*{4. Короткая самопроверка перед решением}
\sectionline
\begin{itemize}
\checkitem{Я понимаю, можно ли считать тело материальной точкой.}
\checkitem{Я выбрал правильную относительную скорость: сумму или модуль разности.}
\checkitem{Я привёл единицы к единой системе.}
\checkitem{В координатной задаче я выбрал ось и поставил знаки скоростей по направлению движения.}
\checkitem{Для встречи я использую именно равенство координат.}
\checkitem{Полученный ответ реалистичен и имеет правильную единицу измерения.}
\end{itemize}

\subsection*{Полезные переводы единиц}
\begin{center}
\begin{tabular}{@{}ll@{}}
\toprule
Известно & Можно быстро получить \\
\midrule
\(1\unit{км/ч}\) & \(\dfrac{5}{18}\unit{м/с}\) \\
\(1\unit{м/с}\) & \(3{,}6\unit{км/ч}\) \\
\(1\unit{мин}\) & \(\dfrac{1}{60}\unit{ч}\) \\
\(1\unit{с}\) & \(\dfrac{1}{3600}\unit{ч}\) \\
\(1\unit{км}\) & \(1000\unit{м}\) \\
\bottomrule
\end{tabular}
\end{center}

\subsection*{Проверка здравого смысла}
Если два объекта движутся навстречу, относительная скорость должна быть больше каждой из отдельных скоростей. При попутном движении она меньше скорости более быстрого объекта. Время полного прохождения увеличивается при увеличении суммарной длины и уменьшается при увеличении относительной скорости.

\clearpage
\section*{5. Тренировочный блок --- Путь А}
\sectionline
\textbf{10 задач.} Решения в этом блоке не приведены. Сложность постепенно возрастает.

\begin{enumerate}
\item Поезд движется со скоростью \(72\unit{км/ч}\) и проходит мимо столба за \(25\unit{с}\). Найдите длину поезда.

\item Поезд движется со скоростью \(90\unit{км/ч}\) и проходит мимо платформы длиной \(650\unit{м}\) за \(50\unit{с}\). Найдите длину поезда.

\item Два поезда длиной \(180\unit{м}\) и \(220\unit{м}\) движутся навстречу друг другу со скоростями \(54\unit{км/ч}\) и \(72\unit{км/ч}\). За какое время они полностью пройдут друг мимо друга?

\item Два поезда длиной \(240\unit{м}\) и \(160\unit{м}\) движутся в одном направлении со скоростями \(90\unit{км/ч}\) и \(54\unit{км/ч}\). Сколько времени потребуется более быстрому поезду, чтобы полностью обогнать более медленный, если в начальный момент передняя точка быстрого поезда находится на уровне задней точки медленного?

\item Два судна длиной \(80\unit{м}\) и \(50\unit{м}\) идут в одном направлении. В начале нос более быстрого судна находится на \(170\unit{м}\) позади кормы медленного. Через \(20\unit{мин}\) после полного обгона корма быстрого судна оказывается на \(300\unit{м}\) впереди носа медленного. На сколько \(\unit{км/ч}\) отличаются их скорости?

\item Из пунктов \(A\) и \(B\), расстояние между которыми \(150\unit{км}\), одновременно навстречу друг другу выехали автомобили со скоростями \(60\unit{км/ч}\) и \(90\unit{км/ч}\). Найдите время и место встречи, считая координату от \(A\).

\item В момент \(t=0\) первое тело находится в точке \(x_1=20\unit{км}\) и движется вправо со скоростью \(15\unit{км/ч}\). Второе находится в точке \(x_2=140\unit{км}\) и движется влево со скоростью \(25\unit{км/ч}\). Найдите время и координату встречи.

\item Первый автомобиль стартует из точки \(x=0\) со скоростью \(72\unit{км/ч}\). Одновременно второй автомобиль находится на \(30\unit{км}\) впереди и движется в том же направлении со скоростью \(54\unit{км/ч}\). Через какое время первый автомобиль догонит второй и на каком расстоянии от начала координат?

\item Тело движется равноускоренно по закону \(x=x_0+v_0t+\dfrac{at^2}{2}\). Найдите координату через \(3\unit{с}\), если \(x_0=5\unit{м}\), \(v_0=4\unit{м/с}\), \(a=2\,\dfrac{\text{м}}{\text{с}^2}\).

\item Поезд длиной \(250\unit{м}\) полностью обгоняет поезд длиной \(150\unit{м}\) за \(40\unit{с}\). Медленный поезд движется со скоростью \(54\unit{км/ч}\). Найдите скорость быстрого поезда.
\end{enumerate}

\clearpage
\section*{6. Ответы}
\sectionline
\renewcommand{\arraystretch}{1.35}
\begin{tabularx}{\linewidth}{@{}>{\bfseries}c X@{}}
\toprule
№ & Ответ \\
\midrule
1 & \(500\unit{м}\). \\
2 & \(600\unit{м}\). \\
3 & \(\dfrac{80}{7}\unit{с}\approx11{,}4\unit{с}\). \\
4 & \(40\unit{с}\). \\
5 & \(1{,}8\unit{км/ч}\). \\
6 & \(t=1\unit{ч}\), место встречи --- \(60\unit{км}\) от \(A\). \\
7 & \(t=3\unit{ч}\), \(x=65\unit{км}\). \\
8 & \(t=\dfrac{5}{3}\unit{ч}=1\unit{ч}\ 40\unit{мин}\), \(x=120\unit{км}\). \\
9 & \(26\unit{м}\). \\
10 & \(90\unit{км/ч}\). \\
\bottomrule
\end{tabularx}

\vfill
{\small\color{textgray}\textbf{Главное:} для протяжённых тел считайте относительное перемещение и относительную скорость; для встречи записывайте координаты и приравнивайте их.}

\end{document}
